\documentclass[11pt]{article}
\usepackage[margin=1in]{geometry}
\usepackage{amsmath,amssymb,amsthm}
\usepackage{graphicx}
\usepackage[hidelinks]{hyperref}

\newtheorem{theorem}{Theorem}
\newtheorem{lemma}{Lemma}
\newtheorem{corollary}{Corollary}
\newcommand{\C}{\mathbb C}
\newcommand{\Log}{\operatorname{Log}}
\newcommand{\B}{\operatorname{B}}
\setlength{\emergencystretch}{2em}

\title{An exact convergence criterion for the Binet factorial expansion}
\author{Candidate full solution to the question in arXiv:2102.04891}
\date{13 August 2026}

\begin{document}
\maketitle

\paragraph{Status.}
Candidate full answer, subject to expert review.  The answer is exact but is
necessarily expressed through an associated ordinary Dirichlet series.  A
germ-bijection below proves that arbitrary Dirichlet coefficient behavior can
occur.  In particular, the original Taylor radius alone cannot decide
convergence.

\section{Source question}

The source is P.~Van Mieghem, \emph{Binet's factorial series and extensions to
Laplace transforms}, arXiv:2102.04891 \cite{source}.  Section~8.4, ``Open
question,'' on source PDF page~38 asks:

\begin{quote}
Given the Taylor series
$f(t)=\sum_{k=0}^{\infty}f_k t^k$ with radius of convergence $R_f$, when does
a factorial series (90) of the Laplace transform $\varphi(z)$ converge?
\end{quote}

\begin{center}
\includegraphics[width=\textwidth]{figures/open_problem_crop.png}
\end{center}

The source fixes $\beta>0$ and $\alpha\in\C$.  With
$w=\beta z+\alpha$, its factorial series (90) can be written
\begin{equation}\label{eq:factorial}
 \mathcal S_{\alpha,\beta}f(z)
 =\beta\sum_{m=0}^{\infty}
   a_m\frac{m!}{\prod_{k=0}^{m}(w+k)}
 =\beta\sum_{m=0}^{\infty}a_m\B(m+1,w),
\end{equation}
where
\begin{equation}\label{eq:G}
 G_{\alpha,\beta}(u)
 =(1-u)^{-\alpha}f\bigl(-\beta\Log(1-u)\bigr)
 =\sum_{m=0}^{\infty}a_m u^m
\end{equation}
as an analytic germ at zero.  Thus $a_m$ is exactly the source coefficient
$\phi_m(\alpha,\beta)$, and it is explicitly computable from the given
Taylor coefficients $f_k$ by the Stirling-number transform in the source.

\section{Proof intuition}

The beta factor in the $m$th factorial-series term is only a smoothly varying
perturbation of a Dirichlet weight:
\[
 \B(m+1,w)=\Gamma(w)m^{-w}\bigl(1+O_w(m^{-1})\bigr).
\]
The perturbing sequence has summable first differences.  Multiplication by
such a sequence preserves even conditional convergence, in both directions.
Consequently the exact convergence question is the ordinary Dirichlet-series
question for the transformed Taylor coefficients $a_m$.

This also yields a sharp radius trichotomy for the series in
\eqref{eq:G}.  A radius above one gives exponential coefficient decay; a
radius below one gives exponential growth along a subsequence and hence
termwise divergence; radius one is precisely the nontrivial Dirichlet case.
Finally, the Binet change of variables is a bijection on analytic germs, so
arbitrary Dirichlet coefficients can arise.  The reduction is therefore
optimal, rather than an artifact of the proof.

\section{Exact criterion}

All powers $m^{-w}$ below use the real logarithm of the positive integer $m$.

\begin{theorem}[Factorial--Dirichlet equivalence]\label{thm:main}
Fix $\beta>0$, $\alpha\in\C$, and a germ $f$ analytic at zero.  Let $a_m$ be
defined by \eqref{eq:G}, let $w=\beta z+\alpha$, and suppose
$\operatorname{Re}w>0$.

Then the factorial series \eqref{eq:factorial} converges if and only if
\begin{equation}\label{eq:dirichlet}
 \sum_{m=1}^{\infty}a_m m^{-w}
\end{equation}
converges.  The same equivalence holds for absolute convergence.  The
$m=0$ term in \eqref{eq:factorial} is the harmless finite term $\beta a_0/w$.
\end{theorem}

\begin{lemma}[Bounded-variation multiplier]\label{lem:bv}
Suppose $c_m\to c\ne0$ and
$\sum_m|c_{m+1}-c_m|<\infty$.  Then $\sum b_m$ converges if and only if
$\sum c_m b_m$ converges.
\end{lemma}

\begin{proof}
If $\sum b_m$ converges, summation by parts applied to its bounded partial
sums and the summable differences of $c_m$ proves convergence of
$\sum c_m b_m$.  Since $c_m$ is bounded away from zero for all sufficiently
large $m$, the sequence $1/c_m$ also converges to $1/c$ and
\[
 \left|\frac1{c_{m+1}}-\frac1{c_m}\right|
 \leq C|c_{m+1}-c_m|.
\]
The same argument with multiplier $1/c_m$ proves the converse.  Finitely many
initial terms do not matter.
\end{proof}

\begin{proof}[Proof of Theorem~\ref{thm:main}]
Euler's beta identity gives
\begin{equation}\label{eq:beta}
 \B(m+1,w)=\frac{\Gamma(m+1)\Gamma(w)}{\Gamma(m+1+w)}
 =\Gamma(w)m^{-w}c_m(w),
 \qquad
 c_m(w)=\frac{m^w\Gamma(m+1)}{\Gamma(m+1+w)}.
\end{equation}
The fixed-parameter gamma-quotient expansion (equivalently, Stirling's
formula with one further term) gives
\begin{equation}\label{eq:quotient}
 c_m(w)=1+O_w(m^{-1}),
 \qquad c_{m+1}(w)-c_m(w)=O_w(m^{-2}).
\end{equation}
Thus $c_m(w)$ satisfies Lemma~\ref{lem:bv}.  Apply the lemma to
$b_m=a_m m^{-w}$ in \eqref{eq:beta}; multiplication by the fixed nonzero
factor $\beta\Gamma(w)$ is irrelevant.  This proves ordinary convergence in
both directions.  Moreover $|c_m(w)|$ is bounded above and away from zero
eventually, which proves the absolute-convergence equivalence.
\end{proof}

\begin{corollary}[Sharp transformed-radius trichotomy]\label{cor:radius}
Let $R_G$ be the radius of convergence of \eqref{eq:G}.
\begin{enumerate}
\item If $R_G>1$, \eqref{eq:factorial} converges absolutely for every
      $z$ with $\operatorname{Re}(\beta z+\alpha)>0$.
\item If $R_G<1$, \eqref{eq:factorial} diverges for every such fixed $z$;
      in fact its terms fail to tend to zero.
\item If $R_G=1$, its convergence and absolute convergence regions are
      exactly those of \eqref{eq:dirichlet}.  Hence they are governed by the
      usual convergence and absolute-convergence abscissae, with the usual
      possibility of exceptional behavior on the boundary lines.
\end{enumerate}
\end{corollary}

\begin{proof}
If $R_G>1$, Cauchy's estimate on a circle of any radius
$1<r<R_G$ gives $|a_m|\leq C_r r^{-m}$, so \eqref{eq:dirichlet} is absolutely
convergent for every $w$.  If $R_G<1$, Cauchy--Hadamard gives
\[
 \limsup_{m\to\infty}|a_m|^{1/m}=R_G^{-1}>1.
\]
The beta factor in \eqref{eq:beta} has $m$th-root limit one, so the terms of
\eqref{eq:factorial} cannot tend to zero.  The final assertion follows
directly from Theorem~\ref{thm:main} and the elementary half-plane theorem for
ordinary Dirichlet series.
\end{proof}

\section{When the sum is the Laplace transform}

The exact convergence statement above concerns the numerical factorial
series.  The following records a clean condition under which it also
represents the intended Laplace transform, without concealing an interchange
of a sum and an improper integral.

\begin{theorem}[Absolute representation]\label{thm:representation}
Suppose the chosen continuation of $f$ along $[0,\infty)$ agrees with
\eqref{eq:G} for $0\leq u<1$.  If $\operatorname{Re}w>0$ and
\eqref{eq:factorial} converges absolutely, then the Laplace integral is
absolutely convergent at $z$ and
\begin{equation}\label{eq:representation}
 \int_0^\infty e^{-zt}f(t)\,dt
 =\beta\sum_{m=0}^{\infty}a_m\B(m+1,w).
\end{equation}
\end{theorem}

\begin{proof}
By Theorem~\ref{thm:main}, absolute convergence is equivalent to
$\sum|a_m|m^{-\operatorname{Re}w}<\infty$.  The positive-real beta factors
are comparable to these powers, so
\[
 \sum_{m=0}^{\infty}|a_m|\B(m+1,\operatorname{Re}w)<\infty.
\]
Since
\[
 \B(m+1,\operatorname{Re}w)
 =\int_0^1 u^m(1-u)^{\operatorname{Re}w-1}\,du,
\]
Tonelli's theorem permits absolute termwise integration.  Binet's substitution
$t=-\beta\log(1-u)$ then gives
\begin{align*}
 \int_0^\infty e^{-zt}f(t)\,dt
 &=\beta\int_0^1(1-u)^{\beta z-1}
       f\bigl(-\beta\log(1-u)\bigr)\,du\\
 &=\beta\int_0^1(1-u)^{w-1}G_{\alpha,\beta}(u)\,du\\
 &=\beta\sum_{m=0}^{\infty}a_m\B(m+1,w).
\end{align*}
The same majorant proves absolute convergence of the first integral.
\end{proof}

No unconditional claim is made here that mere conditional convergence of
\eqref{eq:factorial} justifies \eqref{eq:representation}; that requires an
additional boundary summability hypothesis.

\section{Optimality and a same-radius separation}

\begin{theorem}[Universality of the transformed coefficients]\label{thm:germ}
For fixed $\alpha\in\C$ and $\beta>0$, the transformation
\[
 f(t)\longmapsto
 (1-u)^{-\alpha}f\bigl(-\beta\Log(1-u)\bigr)
\]
is a linear bijection between analytic germs at $t=0$ and analytic germs at
$u=0$.  Its inverse is
\begin{equation}\label{eq:inverse}
 f(t)=e^{-\alpha t/\beta}G\bigl(1-e^{-t/\beta}\bigr).
\end{equation}
Thus arbitrary power-series coefficients $a_m$ can occur in
\eqref{eq:G}; in the unrestricted generality of the source question, the
Dirichlet criterion of Theorem~\ref{thm:main} cannot be replaced by a simpler
universal coefficient phenomenon.
\end{theorem}

\begin{proof}
Both substitutions are analytic near their origins and are inverse there:
$u=1-e^{-t/\beta}$ if and only if $t=-\beta\Log(1-u)$.  Since
$(1-u)^\alpha=e^{-\alpha t/\beta}$ under this substitution, direct composition
proves both inverse identities.
\end{proof}

The next explicit pair shows more concretely that even the complete value of
the original radius $R_f$ carries too little information.

\begin{theorem}[Equal $R_f$, opposite convergence]\label{thm:examples}
Take $\alpha=0$, $\beta=1$, $q=5/2$, and $R=\log q$.  Define for $t\geq0$
\[
 f_A(t)=\frac1{q-e^{-t}},
 \qquad
 f_B(t)=\frac1{1+(t/R)^2}.
\]
Both Laplace transforms converge for $\operatorname{Re}z>0$, and the Taylor
series of both functions at zero has radius exactly $R$.  Nevertheless, the
factorial expansion for $f_A$ converges absolutely for every
$\operatorname{Re}z>0$, whereas that for $f_B$ diverges for every fixed $z$.
\end{theorem}

\begin{proof}
The singularities of $f_A$ are
$-\log q-2\pi i k$, $k\in\mathbb Z$, so its nearest singularity is $-R$.
The nearest singularities of $f_B$ are $\pm iR$.  Both Taylor radii are
therefore $R$.  On the positive real axis, $f_A$ is bounded and $f_B(t)$ is
$O(t^{-2})$, so both Laplace integrals converge when
$\operatorname{Re}z>0$.

Their transformed germs are
\[
 G_A(u)=f_A(-\Log(1-u))=\frac1{q-1+u},
\]
whose radius is $q-1=3/2>1$, and
\[
 G_B(u)=\frac1{1+\bigl(-\Log(1-u)/R\bigr)^2}.
\]
The latter has poles at $u=1-e^{\mp iR}$, of modulus
$2\sin(R/2)$.  Since
\[
 e^{\pi/3}>1+\frac\pi3+\frac12\left(\frac\pi3\right)^2>\frac52,
\]
we have $R=\log(5/2)<\pi/3$, and hence
$2\sin(R/2)<2\sin(\pi/6)=1$.  Thus $G_B$ has radius below one.
Corollary~\ref{cor:radius} gives the two opposite conclusions.
\end{proof}

\section{Scope, verification, and novelty}

The theorem answers convergence at every parameter point in the source
half-plane $\operatorname{Re}(\beta z+\alpha)>0$.  It also gives an absolute
criterion ensuring equality with the Laplace integral.  It deliberately does
not infer equality from conditional convergence alone.

The gamma multiplier, both branches of the germ inverse, the two example
radii, and the Tonelli step were audited independently in the accompanying
verification report.  No numerical computation is used as proof.

The four lightweight run indexes had no entry for this arXiv id or problem.
A bounded primary-arXiv search on 13 August 2026 used the exact title and close
phrases involving factorial series, Laplace transforms, Taylor coefficients,
and convergence; it found the source but no later solution.  The source itself
already points out a classical $\beta=1$ comparison with a Dirichlet series,
citing Lauwerier and Temme.  That comparison is not claimed as new.  The
candidate contribution is the rigorous all-parameter equivalence including
conditional convergence, the sharp transformed-radius trichotomy, the
germ-universality/optimality theorem, and the equal-$R_f$ separation.  Novelty
confidence is therefore moderate rather than high pending a specialist check
of the classical factorial-series literature.

\paragraph{Human-review recommendation.}
Review as a candidate full answer.  The highest-value checks are
Theorem~\ref{thm:main}'s two-way conditional-convergence argument and whether
an equivalent universality/radius formulation already occurs in the older
literature cited by the source.

\begin{thebibliography}{9}
\bibitem{source}
\raggedright P. Van Mieghem,
\emph{Binet's factorial series and extensions to Laplace transforms},
arXiv:2102.04891v7 (2023).\par
\end{thebibliography}

\end{document}
